\section*{第 \ref{ch:b1:eukaryotic-cell} 章——真核细胞的功能组织}

\begin{solution}{exo:b1:eukaryotic-cell:1}
\omterm{def:b1:eukaryotic-cell:nucleus}{细胞核}（染色体、转录）；\omterm{def:b1:eukaryotic-cell:endomembrane}{粗面内质网}（分泌性蛋白质和膜蛋白）；
\omterm{def:b1:eukaryotic-cell:endomembrane}{滑面内质网}（脂质、钙、解毒）；\omterm{def:b1:eukaryotic-cell:endomembrane}{高尔基体}（分类、修饰、包装）；
\omterm{def:b1:eukaryotic-cell:endomembrane}{溶酶体}（消化）；\omterm{def:b1:eukaryotic-cell:peroxisome}{过氧化物酶体}（氧化反应、过氧化氢酶）；
\omterm{def:b1:eukaryotic-cell:mitochondrion}{线粒体}（呼吸作用、ATP）。
\end{solution}

\begin{solution}{exo:b1:eukaryotic-cell:2}
由结合在\omterm{def:b1:eukaryotic-cell:endomembrane}{粗面内质网}上的核糖体合成，并穿入其腔内（\omterm{prop:b1:eukaryotic-cell:secretory}{信号序列}）；
囊泡从\omterm{def:b1:eukaryotic-cell:endomembrane}{内质网}出芽，并与\omterm{def:b1:eukaryotic-cell:endomembrane}{高尔基体}顺面融合；
蛋白质穿过\omterm{def:b1:eukaryotic-cell:endomembrane}{高尔基体}堆叠结构，被修饰并分类；
随后从反面离开，进入正在成熟为分泌颗粒的浓缩泡；
颗粒在信号作用下与顶端膜融合，
并将内容物排入\omterm{def:b1:flowering-plant-organization:tissues}{导管}。
\end{solution}

\begin{solution}{exo:b1:eukaryotic-cell:3}
$7 + 32 = \qty{39}{\%}$。内膜向内折叠成嵴，
因为其上载有呼吸链和 ATP 合酶，二者的通量与膜面积成正比；
外膜则是简单的包膜。
\end{solution}

\begin{solution}{exo:b1:eukaryotic-cell:4}
植物：纤维素\omterm{def:b1:eukaryotic-cell:plantcell}{细胞壁}、大\omterm{def:b1:eukaryotic-cell:plantcell}{液泡}、\omterm{def:b1:eukaryotic-cell:plastid}{质体}（另有\omterm{def:b1:eukaryotic-cell:plantcell}{胞间连丝}）。
动物：\omterm{def:b1:eukaryotic-cell:endomembrane}{溶酶体}、中心粒（以及
爬行和吞噬的能力）。
\end{solution}

\begin{solution}{exo:b1:eukaryotic-cell:5}
\omterm{def:b1:eukaryotic-cell:endomembrane}{内质网}在 \qty{3}{min}（已开始下降），\omterm{def:b1:eukaryotic-cell:endomembrane}{高尔基体}约在 \qty{15}{min}，
颗粒约在 \qty{60}{min}。从\omterm{def:b1:eukaryotic-cell:endomembrane}{内质网}到颗粒的转运约
需 \qty{45}{min}；约在 \qty{110}{min} 时分泌一半。
\end{solution}

\begin{solution}{exo:b1:eukaryotic-cell:6}
\omterm{def:b1:eukaryotic-cell:endomembrane}{高尔基体}借\omterm{def:b1:eukaryotic-cell:cytoskeleton}{微管}上的动力蛋白固定在\omterm{def:b1:eukaryotic-cell:cytoskeleton}{中心体}附近，
碎裂并分散；囊泡运输减慢为扩散，向外周的运输停止；
纺锤体无法形成，\omterm{def:b1:cell-unit-of-life:cell}{细胞}停滞在有丝分裂期；
纤毛的核心是\omterm{def:b1:eukaryotic-cell:cytoskeleton}{微管}，其形成随之停止，
已有纤毛较稳定，但不再更新。
\end{solution}

\begin{solution}{exo:b1:eukaryotic-cell:7}
具有双层膜，内层为细菌型；环状 DNA，不含组蛋白；
核糖体大小与细菌相似，对抗生素敏感；以裂变方式分裂；
基因序列与α-变形菌聚为一类。它无法独立生存，
因为其大部分基因已转移到\omterm{def:b1:eukaryotic-cell:nucleus}{细胞核}中，它需从
\omterm{def:b1:eukaryotic-cell:organelle}{细胞溶胶}中输入大部分蛋白质。
\end{solution}

\begin{solution}{exo:b1:eukaryotic-cell:8}
糖类在\omterm{def:b1:eukaryotic-cell:endomembrane}{内质网}和\omterm{def:b1:eukaryotic-cell:endomembrane}{高尔基体}的腔面一侧加上。每次囊泡融合都保持
膜的方向性：囊泡的腔面成为下一区室的腔面，
囊泡与\omterm{def:b1:cell-unit-of-life:cell}{质膜}融合时，其内腔向外敞开，因此
糖类朝向外部。
就拓扑而言，腔面即外侧。
\end{solution}

\begin{solution}{exo:b1:eukaryotic-cell:9}
每个颗粒 $\pi d^2 = \qty{3.14}{\micro m^2}$；每天
\qty{6280}{\micro m^2}，即 \qty{157}{} 倍于顶端面。
\omterm{def:b1:cell-unit-of-life:cell}{细胞}必须以同样速率通过\omterm{def:b1:eukaryotic-cell:endocytosis}{胞吞作用}回收膜，因此
顶端膜每九分钟更替一次，被回收的
膜被送回\omterm{def:b1:eukaryotic-cell:endomembrane}{高尔基体}再利用。
\end{solution}

\begin{solution}{exo:b1:eukaryotic-cell:10}
\qty{320}{} 个残基的肽链是前体，其\omterm{prop:b1:eukaryotic-cell:secretory}{信号序列}为 \qty{20}{} 个残基，
\omterm{prop:b1:eukaryotic-cell:secretory}{信号序列}通常在进入\omterm{def:b1:eukaryotic-cell:endomembrane}{内质网}时被切除。
有囊泡时，肽链在合成过程中进入腔内，
\omterm{prop:b1:eukaryotic-cell:secretory}{信号序列}被切除，膜保护产物免受蛋白酶的作用。
缺失\omterm{prop:b1:eukaryotic-cell:secretory}{信号序列}的突变体不被\omterm{def:b1:eukaryotic-cell:endomembrane}{内质网}的机制识别：
它在游离核糖体上合成，停留在\omterm{def:b1:eukaryotic-cell:organelle}{细胞溶胶}中，全长，
从不被分泌。
\end{solution}

\begin{solution}{exo:b1:eukaryotic-cell:11}
通过\omterm{def:b1:eukaryotic-cell:endocytosis}{胞吞作用}内化的膜和老化的\omterm{def:b1:eukaryotic-cell:organelle}{细胞器}被运送到\omterm{def:b1:eukaryotic-cell:endomembrane}{溶酶体}中消化；
若缺乏该酶，其所含的脂质无法分解，
便会充满\omterm{def:b1:eukaryotic-cell:endomembrane}{溶酶体}，使\omterm{def:b1:eukaryotic-cell:endomembrane}{溶酶体}肿胀，
直至\omterm{def:b1:cell-unit-of-life:cell}{细胞}衰竭。膜更替最快的\omterm{def:b1:cell-unit-of-life:cell}{细胞}
（\omterm{def:b1:body-plans-tissues:nervous}{神经元}持续回收突触囊泡；巨噬\omterm{def:b1:cell-unit-of-life:cell}{细胞}）积累得最快。
注入的酶若携带\omterm{def:b1:eukaryotic-cell:endomembrane}{溶酶体}酶所具有的糖类标记，
便能与\omterm{def:b1:cell-unit-of-life:cell}{细胞}表面的受体结合，
并经受体介导的\omterm{def:b1:eukaryotic-cell:endocytosis}{胞吞作用}进入\omterm{def:b1:eukaryotic-cell:endomembrane}{溶酶体}——
这正是酶替代疗法的原理。
\end{solution}

\begin{solution}{exo:b1:eukaryotic-cell:12}
内共生学说能很好地解释\omterm{def:b1:eukaryotic-cell:mitochondrion}{线粒体}和\omterm{def:b1:eukaryotic-cell:plastid}{质体}（二者均有双层膜、
拥有自身的基因组和核糖体，以裂变方式分裂，且有细菌类的近缘种），
因而也能解释真核生物的能量代谢。它无法解释\omterm{def:b1:eukaryotic-cell:nucleus}{细胞核}、\omterm{def:b1:eukaryotic-cell:endomembrane}{内质网}、
\omterm{def:b1:eukaryotic-cell:endomembrane}{高尔基体}和\omterm{def:b1:eukaryotic-cell:cytoskeleton}{细胞骨架}，这些结构在细菌中没有对应物，
似乎是宿主谱系的新发明——它们是其膜的内陷和新的蛋白质。
“联合体”一词对各\omterm{def:b1:eukaryotic-cell:organelle}{细胞器}较为贴切；
“由膜系统维系在一起”则指向较大的、尚未得到解释的部分：
即一个在吞入细菌之前，
结构上已经是真核的宿主\omterm{def:b1:cell-unit-of-life:cell}{细胞}。
\end{solution}

\begin{solution}{pb:b1:eukaryotic-cell:1}
\textbf{1.} 短时脉冲标记一小批分子，使其位置标示出某一时刻；
过量的冷氨基酸阻止进一步掺入，
从而使这批分子不会因后续合成而模糊不清。
\textbf{2.} \omterm{def:b1:eukaryotic-cell:endomembrane}{内质网}：\qty{3}{min}（最早时点）；\omterm{def:b1:eukaryotic-cell:endomembrane}{高尔基体}与浓缩
泡：约 \qty{15}{min}；颗粒：约 \qty{60}{min}。
\textbf{3.} 从\omterm{def:b1:eukaryotic-cell:endomembrane}{内质网}到\omterm{def:b1:eukaryotic-cell:endomembrane}{高尔基体}约 \qty{10}{min}；从\omterm{def:b1:eukaryotic-cell:endomembrane}{高尔基体}到颗粒约
\qty{45}{min}。
\textbf{4.} 大部分位于颗粒中（\qty{70}{\%}），十分之一在\omterm{def:b1:eukaryotic-cell:endomembrane}{高尔基体}中，
十分之一已分泌：组成型途径和最初形成的颗粒释放了一部分，
而大多数颗粒则等待刺激。
\textbf{5.} 囊泡出芽和运输需要能量以及流动的膜；
低温下，从\omterm{def:b1:eukaryotic-cell:endomembrane}{内质网}到\omterm{def:b1:eukaryotic-cell:endomembrane}{高尔基体}的步骤停止，
蛋白质在阻断处的上游积累。
\textbf{6.} $10\,\mathrm{g}/\num{8e10} = \qty{1.25e-10}{g} =
\qty{125}{pg}$ 每个\omterm{def:b1:cell-unit-of-life:cell}{细胞}每天。
\textbf{7.} $V = \frac{\pi}{6}(1)^3 = \qty{0.52}{\micro m^3} =
\qty{5.2e-13}{mL}$；蛋白质 $\qty{1.05e-13}{g} = \qty{0.1}{pg}$。
\textbf{8.} $125/0.1 = 1250$ 个颗粒每天，每分钟 $0.87$ 个。
\textbf{9.} $V = 12^3 = \qty{1728}{\micro m^3} = \qty{1.73e-9}{cm^3}$;
质量为 \qty{1.8e-9}{g} = \qty{1800}{pg}。每天分泌
其质量的 $125/1800 = 7\%$（整个\omterm{def:b1:cell-unit-of-life:cell}{细胞}仅 \qty{20}{\%} 为蛋白质，故
约占其蛋白质含量的三分之一）。
\textbf{10.} 摩尔质量 $250\times 110 = \qty{27500}{g/mol}$；
$\qty{1.05e-13}{g}/27500\times\num{6e23} = \num{2.3e6}$ 个分子。
\textbf{11.} $250/5 = \qty{50}{s}$.
\textbf{12.} 每天的分子数 $1250\times\num{2.3e6} = \num{2.9e9}$；
一个核糖体每天制造 $86\,400/50 = 1728$ 个；共需
$\num{1.7e6}$ 个核糖体，约占\omterm{def:b1:cell-unit-of-life:cell}{细胞}全部核糖体的 \qty{40}{\%}。
\textbf{13.} $\pi d^2 = \qty{3.14}{\micro m^2}$.
\textbf{14.} $1250\times 3.14 = \qty{3900}{\micro m^2}$.
\textbf{15.} $3900/144 = 27$ 倍：\omterm{def:b1:cell-unit-of-life:cell}{细胞}通过
\omterm{def:b1:eukaryotic-cell:endocytosis}{胞吞作用}回收相同面积的膜，并将其送回\omterm{def:b1:eukaryotic-cell:endomembrane}{高尔基体}。
\textbf{16.} $20\,000/3900 = 5$ 天：若无回收，\omterm{def:b1:eukaryotic-cell:endomembrane}{内质网}只能供应
五天的颗粒膜；膜必须
持续循环。
\textbf{17.} 每根微绒毛 $\pi\times 0.1\times 1 =
\qty{0.31}{\micro m^2}$；2000 根合计 \qty{630}{\micro m^2}；倍数为 $5.4$。
\textbf{18.} 位于肽链起始端的是\omterm{prop:b1:eukaryotic-cell:secretory}{信号序列}，
由二十个主要为疏水的残基组成。
\textbf{19.} \omterm{prop:b1:eukaryotic-cell:secretory}{信号序列}引导生长中的肽链进入囊泡，
信号肽酶将其切除；膜保护蛋白质
免受蛋白酶的作用。
\textbf{20.} 只有带\omterm{prop:b1:eukaryotic-cell:secretory}{信号序列}的蛋白质才能进入\omterm{def:b1:eukaryotic-cell:endomembrane}{内质网}；\omterm{def:b1:eukaryotic-cell:organelle}{细胞溶胶}中的
蛋白质缺少\omterm{prop:b1:eukaryotic-cell:secretory}{信号序列}，留在外部。
\textbf{21.} 该酶在游离核糖体上合成，并留在
\omterm{def:b1:eukaryotic-cell:organelle}{细胞溶胶}中；作为消化酶（以无活性前体存在时危害不大，
但一旦被激活）会从内部消化\omterm{def:b1:cell-unit-of-life:cell}{细胞}。
\textbf{22.} 在\omterm{def:b1:eukaryotic-cell:endomembrane}{高尔基体}反面：\omterm{def:b1:eukaryotic-cell:endomembrane}{溶酶体}酶携带一种
可被受体识别的标记（磷酸化的糖），
由该受体将其包装进运往\omterm{def:b1:eukaryotic-cell:endomembrane}{溶酶体}的囊泡，而非
分泌颗粒。
\textbf{23.} 在\omterm{def:b1:eukaryotic-cell:endomembrane}{内质网}（及\omterm{def:b1:eukaryotic-cell:endomembrane}{高尔基体}）中：囊泡形成、
沿\omterm{def:b1:eukaryotic-cell:cytoskeleton}{微管}的马达蛋白运输和膜融合均消耗 ATP；同时
蛋白质合成本身也会停止，因此脉冲甚至无法完成。
\textbf{24.} 按每分钟 $0.87$ 个颗粒计，\qty{45}{min} 的转运
途中蛋白质约相当于 \qty{40}{} 个颗粒的量。
\textbf{25.} 从合成到颗粒的转运约 \qty{45}{min}
（若计入\omterm{def:b1:eukaryotic-cell:endomembrane}{内质网}则为 \qty{55}{min}），在实验中约两小时后分泌，
（在活体腺体中为四小时或更久，因为颗粒
要等待进食后才释放）；每分钟释放 \qty{0.9}{} 个颗粒，
每天分泌 \qty{7}{\%} 的\omterm{def:b1:cell-unit-of-life:cell}{细胞质}量。
\end{solution}
